% GENERATED FILE. Edit canonical lessons, then run npm run generate:books.
% canonical-source-hash: fnv1a64:6f8e52650ddd2583
% canonical-lessons: PA-C211-verbs-return, PA-C211-janda, PA-C211-janda-review, PA-C211-khanda-pinda, PA-C211-khanda-review, PA-C211-tea-every-day, PA-C211-payoff

\chapter{Going and Drinking Every Day}
\label{ch:pa-a2-vowel-stem}
\hlchaptermodality{211}

\begin{chapteropening}
\textbf{By the end of this chapter:} \emph{I can say what I go to, eat and drink every day, with the present forms of jana, khana and pina.}

Say where you go and what you eat and drink every day, as a man, a woman and a group.
\end{chapteropening}

\section[\texorpdfstring{\pa{ਜਾਣਾ} · \pa{ਖਾਣਾ} … (jāṇā · khāṇā …)}{ਜਾਣਾ · ਖਾਣਾ … (jāṇā · khāṇā …)}]{Return to four familiar verbs}

\label{lesson:PA-C211-verbs-return}



\begin{quote}
\emph{Hear:} \emph{jāṇā; khāṇā; pīṇā} Say \textbf{\pa{ਜਾਣਾ} · \pa{ਖਾਣਾ} · \pa{ਪੀਣਾ}}.
\end{quote}

\begin{grammarlens}[title={Grammar Lens}]
You already say \textbf{\pa{ਕਰਦਾ}}, \textquotedblleft{}does\textquotedblright{}, and \textbf{\pa{ਸੌਂਦਾ}}, \textquotedblleft{}sleeps\textquotedblright{}. Listen to the two middles: \emph{kar-dā} runs straight into its ending, but \emph{saun-dā} carries a nasal hum before it. The next lesson gives that hum to three verbs you use every day.
\end{grammarlens}

\subsection*{Your turn}
\begin{itemize}
\raggedright
  \item \emph{Hear:} \emph{to go} Say \textbf{\pa{ਜਾਣਾ}}.
  \item \emph{Hear:} \emph{to eat} Say \textbf{\pa{ਖਾਣਾ}}.
  \item \emph{Hear:} \emph{to drink} Say \textbf{\pa{ਪੀਣਾ}}.
  \item \emph{Hear:} \emph{sleeps, of a man} Say \textbf{\pa{ਸੌਂਦਾ}}.
\end{itemize}

\begin{tcolorbox}[breakable,colback=teal!4,colframe=teal!35!black,title={Before you move on}]
Say the four once more, slowly. Nothing new is needed yet.
\end{tcolorbox}
\section[\texorpdfstring{\pa{ਜਾਂਦਾ} · \pa{ਜਾਂਦੀ} (jāndā · jāndī)}{ਜਾਂਦਾ · ਜਾਂਦੀ (jāndā · jāndī)}]{\pa{ਜਾਂਦਾ}, \pa{ਜਾਂਦੀ} — goes, every day}

\label{lesson:PA-C211-janda}



\begin{quote}
Say \textbf{\pa{ਜਾਣਾ}}, \textquotedblleft{}to go\textquotedblright{}, then \textbf{\pa{ਸੌਂਦਾ}}, \textquotedblleft{}sleeps\textquotedblright{}. Then \emph{every day}: \textbf{\pa{ਰੋਜ਼}}.
\end{quote}

\begin{grammarlens}[title={Grammar Lens: a stem that ends in a vowel}]
Take \textbf{-\pa{ਣਾ}} off \textbf{\pa{ਜਾਣਾ}} and the stem \textbf{\pa{ਜਾ}} ends in a vowel. Before \textbf{-\pa{ਦਾ}} such a stem takes a nasal, the dot \textbf{\pa{ਂ}}, as \textbf{\pa{ਸੌਂਦਾ}} does:

\begin{quote}
\textbf{\pa{ਜਾ}} + \textbf{\pa{ਂ}} + \textbf{\pa{ਦਾ}} $\to$ \textbf{\pa{ਜਾਂਦਾ}}
\end{quote}

\begin{quote}
\textbf{\pa{ਮੈਂ} \pa{ਰੋਜ਼} \pa{ਸਕੂਲ} \pa{ਜਾਂਦਾ} \pa{ਹਾਂ।}} — \emph{maiṁ roz sakūl jāndā hāṁ.} — I go to school every day (a man speaking).
\end{quote}

\begin{quote}
\textbf{\pa{ਮੈਂ} \pa{ਰੋਜ਼} \pa{ਸਕੂਲ} \pa{ਜਾਂਦੀ} \pa{ਹਾਂ।}} — \emph{maiṁ roz sakūl jāndī hāṁ.} — I go to school every day (a woman speaking).
\end{quote}

The ending agrees with the one who goes, as in \textbf{\pa{ਮੈਂ} \pa{ਪੰਜਾਬੀ} \pa{ਬੋਲਦਾ} \pa{ਹਾਂ}}. The helper \textbf{\pa{ਹਾਂ}} stays at the end.
\end{grammarlens}

\subsection*{Your turn}
\begin{itemize}
\raggedright
  \item \emph{Hear:} \emph{I go home every day, as a man} Say \textbf{\pa{ਮੈਂ} \pa{ਰੋਜ਼} \pa{ਘਰ} \pa{ਜਾਂਦਾ} \pa{ਹਾਂ।}}
  \item \emph{Hear:} \emph{she goes to school every day} Say \textbf{\pa{ਉਹ} \pa{ਰੋਜ਼} \pa{ਸਕੂਲ} \pa{ਜਾਂਦੀ} \pa{ਹੈ।}}
  \item \emph{Read it:} \textbf{\pa{ਜਾਂਦਾ}} Point to the dot that brings the nasal.
  \item \emph{Write it:} copy \textbf{\pa{ਜਾਂਦਾ}} once
\end{itemize}

\begin{tcolorbox}[breakable,colback=teal!4,colframe=teal!35!black,title={Before you move on}]
\emph{Goes}, of a man? (\textbf{\pa{ਜਾਂਦਾ}}.) Of a woman? (\textbf{\pa{ਜਾਂਦੀ}}.)
\end{tcolorbox}
\section[\texorpdfstring{\pa{ਜਾਂਦਾ} · \pa{ਜਾਂਦੀ} (jāndā · jāndī)}{ਜਾਂਦਾ · ਜਾਂਦੀ (jāndā · jāndī)}]{Hear and say goes again}

\label{lesson:PA-C211-janda-review}



\begin{quote}
\emph{Hear:} \emph{jāndā; jāndī} Say \textbf{\pa{ਜਾਂਦਾ} · \pa{ਜਾਂਦੀ}}.
\end{quote}

\subsection*{Your turn}
\begin{itemize}
\raggedright
  \item \emph{Hear:} \emph{he goes home every day} Say \textbf{\pa{ਉਹ} \pa{ਰੋਜ਼} \pa{ਘਰ} \pa{ਜਾਂਦਾ} \pa{ਹੈ।}}
  \item \emph{Hear:} \emph{I go to school every day, as a woman} Say \textbf{\pa{ਮੈਂ} \pa{ਰੋਜ਼} \pa{ਸਕੂਲ} \pa{ਜਾਂਦੀ} \pa{ਹਾਂ।}}
  \item \emph{Read it:} \textbf{\pa{ਜਾਂਦੀ}} Say \emph{jāndī}.
  \item \emph{Write it:} cover the model and write \textbf{\pa{ਜਾਂਦੀ}} once
\end{itemize}

\begin{tcolorbox}[breakable,colback=teal!4,colframe=teal!35!black,title={Before you move on}]
Say one true sentence about where you go every day.
\end{tcolorbox}
\section[\texorpdfstring{\pa{ਖਾਂਦਾ} · \pa{ਪੀਂਦਾ} (khāndā · pīndā)}{ਖਾਂਦਾ · ਪੀਂਦਾ (khāndā · pīndā)}]{\pa{ਖਾਂਦਾ}, \pa{ਪੀਂਦਾ} — eats, drinks}

\label{lesson:PA-C211-khanda-pinda}



\begin{quote}
Say \emph{I go to school every day}. Then \textbf{\pa{ਖਾਣਾ}} and \textbf{\pa{ਪੀਣਾ}}.
\end{quote}

\begin{grammarlens}[title={Grammar Lens: \pa{ਖਾ} and \pa{ਪੀ}}]
\textbf{\pa{ਖਾ}} and \textbf{\pa{ਪੀ}} end in vowels too, so they work like \textbf{\pa{ਜਾ}}: \textbf{\pa{ਖਾਂਦਾ}}, \textbf{\pa{ਖਾਂਦੀ}}; \textbf{\pa{ਪੀਂਦਾ}}, \textbf{\pa{ਪੀਂਦੀ}}.

\begin{quote}
\textbf{\pa{ਮੈਂ} \pa{ਰੋਜ਼} \pa{ਰੋਟੀ} \pa{ਖਾਂਦਾ} \pa{ਹਾਂ।}} — \emph{maiṁ roz roṭī khāndā hāṁ.} — I eat roti every day (a man speaking).
\end{quote}

\begin{quote}
\textbf{\pa{ਉਹ} \pa{ਰੋਜ਼} \pa{ਦੁੱਧ} \pa{ਪੀਂਦੀ} \pa{ਹੈ।}} — \emph{uh roz duddh pīndī hai.} — She drinks milk every day.
\end{quote}
\end{grammarlens}

\subsection*{Your turn}
\begin{itemize}
\raggedright
  \item \emph{Hear:} \emph{I drink water every day, as a woman} Say \textbf{\pa{ਮੈਂ} \pa{ਰੋਜ਼} \pa{ਪਾਣੀ} \pa{ਪੀਂਦੀ} \pa{ਹਾਂ।}}
  \item \emph{Hear:} \emph{he eats roti every day} Say \textbf{\pa{ਉਹ} \pa{ਰੋਜ਼} \pa{ਰੋਟੀ} \pa{ਖਾਂਦਾ} \pa{ਹੈ।}}
  \item \emph{Read it:} \textbf{\pa{ਪੀਂਦਾ}} Point to \textbf{\pa{ਂ}}, then say \emph{pīndā}.
  \item \emph{Write it:} copy \textbf{\pa{ਖਾਂਦਾ}} once
\end{itemize}

\begin{tcolorbox}[breakable,colback=teal!4,colframe=teal!35!black,title={Before you move on}]
\emph{Eats}, of a man? (\textbf{\pa{ਖਾਂਦਾ}}.) \emph{Drinks}, of a woman? (\textbf{\pa{ਪੀਂਦੀ}}.)
\end{tcolorbox}
\section[\texorpdfstring{\pa{ਖਾਂਦਾ} · \pa{ਪੀਂਦਾ} (khāndā · pīndā)}{ਖਾਂਦਾ · ਪੀਂਦਾ (khāndā · pīndā)}]{Hear and say eats and drinks again}

\label{lesson:PA-C211-khanda-review}



\begin{quote}
\emph{Hear:} \emph{khāndā; pīndā; jāndā} Say \textbf{\pa{ਖਾਂਦਾ} · \pa{ਪੀਂਦਾ} · \pa{ਜਾਂਦਾ}}.
\end{quote}

\subsection*{Your turn}
\begin{itemize}
\raggedright
  \item \emph{Hear:} \emph{I drink tea every day, as a man} Say \textbf{\pa{ਮੈਂ} \pa{ਰੋਜ਼} \pa{ਚਾਹ} \pa{ਪੀਂਦਾ} \pa{ਹਾਂ।}}
  \item \emph{Hear:} \emph{she eats roti every day} Say \textbf{\pa{ਉਹ} \pa{ਰੋਜ਼} \pa{ਰੋਟੀ} \pa{ਖਾਂਦੀ} \pa{ਹੈ।}}
  \item \emph{Read it:} \textbf{\pa{ਖਾਂਦੀ}} Say \emph{khāndī}.
  \item \emph{Write it:} cover the model and write \textbf{\pa{ਪੀਂਦੀ}} once
\end{itemize}

\begin{tcolorbox}[breakable,colback=teal!4,colframe=teal!35!black,title={Before you move on}]
Say what you drink every day, with the ending that fits you.
\end{tcolorbox}
\section[\texorpdfstring{\pa{ਅਸੀਂ} \pa{ਰੋਜ਼} \pa{ਚਾਹ} \pa{ਪੀਂਦੇ} \pa{ਹਾਂ।} (asīṁ roz chāh …)}{ਅਸੀਂ ਰੋਜ਼ ਚਾਹ ਪੀਂਦੇ ਹਾਂ। (asīṁ roz chāh …)}]{\pa{ਅਸੀਂ} \pa{ਰੋਜ਼} \pa{ਚਾਹ} \pa{ਪੀਂਦੇ} \pa{ਹਾਂ} — we drink tea every day}

\label{lesson:PA-C211-tea-every-day}



\begin{quote}
Say \emph{I drink tea every day}. Then the familiar \textbf{\pa{ਅਸੀਂ} \pa{ਮਿਲਦੇ} \pa{ਹਾਂ}}, \textquotedblleft{}we meet\textquotedblright{}.
\end{quote}

\begin{grammarlens}[title={Grammar Lens: -\pa{ਦੇ} for more than one}]
\textbf{\pa{ਮਿਲਦੇ}} showed the \textbf{-\pa{ਦੇ}} ending for \textbf{\pa{ਅਸੀਂ}}. The vowel stems take it too, nasal and all: \textbf{\pa{ਜਾਂਦੇ}}, \textbf{\pa{ਖਾਂਦੇ}}, \textbf{\pa{ਪੀਂਦੇ}}. \textbf{\pa{ਤੁਸੀਂ}} takes it as well, with \textbf{\pa{ਹੋ}}:

\begin{quote}
\textbf{\pa{ਅਸੀਂ} \pa{ਰੋਜ਼} \pa{ਚਾਹ} \pa{ਪੀਂਦੇ} \pa{ਹਾਂ।}} — \emph{asīṁ roz chāh pīnde hāṁ.} — We drink tea every day.
\end{quote}

\begin{quote}
\textbf{\pa{ਤੁਸੀਂ} \pa{ਰੋਜ਼} \pa{ਸਕੂਲ} \pa{ਜਾਂਦੇ} \pa{ਹੋ।}} — \emph{tusīṁ roz sakūl jānde ho.} — You go to school every day.
\end{quote}
\end{grammarlens}

\subsection*{Your turn}
\begin{itemize}
\raggedright
  \item \emph{Hear:} \emph{we eat roti every day} Say \textbf{\pa{ਅਸੀਂ} \pa{ਰੋਜ਼} \pa{ਰੋਟੀ} \pa{ਖਾਂਦੇ} \pa{ਹਾਂ।}}
  \item \emph{Hear:} \emph{you go home every day} Say \textbf{\pa{ਤੁਸੀਂ} \pa{ਰੋਜ਼} \pa{ਘਰ} \pa{ਜਾਂਦੇ} \pa{ਹੋ।}}
  \item \emph{Read it:} \textbf{\pa{ਅਸੀਂ} \pa{ਰੋਜ਼} \pa{ਚਾਹ} \pa{ਪੀਂਦੇ} \pa{ਹਾਂ।}}
  \item \emph{Write it:} copy \textbf{\pa{ਪੀਂਦੇ}} once
\end{itemize}

\begin{tcolorbox}[breakable,colback=teal!4,colframe=teal!35!black,title={Before you move on}]
\emph{We go}? (\textbf{\pa{ਅਸੀਂ} \pa{ਜਾਂਦੇ} \pa{ਹਾਂ}}.) \emph{You drink}? (\textbf{\pa{ਤੁਸੀਂ} \pa{ਪੀਂਦੇ} \pa{ਹੋ}}.)
\end{tcolorbox}
\section[\texorpdfstring{\pa{ਮੈਂ} \pa{ਰੋਜ਼} \pa{ਚਾਹ} \pa{ਪੀਂਦਾ} \pa{ਹਾਂ।} (maiṁ roz chāh …)}{ਮੈਂ ਰੋਜ਼ ਚਾਹ ਪੀਂਦਾ ਹਾਂ। (maiṁ roz chāh …)}]{Every day: a four-skill payoff}

\label{lesson:PA-C211-payoff}



\begin{quote}
Say \emph{goes}, \emph{eats} and \emph{drinks}, for a man and then for a woman.
\end{quote}

\subsection*{Your turn: four skills}
\begin{itemize}
\raggedright
  \item \emph{Hear:} \emph{I drink tea every day} Say it as yourself: \textbf{\pa{ਮੈਂ} \pa{ਰੋਜ਼} \pa{ਚਾਹ} \pa{ਪੀਂਦਾ} \pa{ਹਾਂ।}} or \textbf{\pa{ਮੈਂ} \pa{ਰੋਜ਼} \pa{ਚਾਹ} \pa{ਪੀਂਦੀ} \pa{ਹਾਂ।}}
  \item \emph{Hear:} \emph{we go to school every day} Say \textbf{\pa{ਅਸੀਂ} \pa{ਰੋਜ਼} \pa{ਸਕੂਲ} \pa{ਜਾਂਦੇ} \pa{ਹਾਂ।}}
  \item \emph{Read it:} \textbf{\pa{ਉਹ} \pa{ਰੋਜ਼} \pa{ਰੋਟੀ} \pa{ਖਾਂਦੀ} \pa{ਹੈ।}} Give its meaning: she eats roti every day.
  \item \emph{Write it:} cover the line and write \textbf{\pa{ਮੈਂ} \pa{ਰੋਜ਼} \pa{ਪਾਣੀ} \pa{ਪੀਂਦਾ} \pa{ਹਾਂ।}} or its woman's form
\end{itemize}

\begin{tcolorbox}[breakable,colback=teal!4,colframe=teal!35!black,title={Before you move on}]
Tell someone three things you do every day: where you go, what you eat, what you drink.
\end{tcolorbox}
